%%%%%%%%%%%%%%%%%%%%%%%%%%%%%%%%%%%%%%%%%%%%%%%%%%%%%%%%%%%%%%%%%%%%%%%%%%%%%%%%
% Official RA-L initial-submission style.
% Source template: RA-L author page -> PaperCept ieeeconf.zip.
% RA-L initial submissions use:
% \documentclass[letterpaper, 10 pt, conference]{ieeeconf}
%%%%%%%%%%%%%%%%%%%%%%%%%%%%%%%%%%%%%%%%%%%%%%%%%%%%%%%%%%%%%%%%%%%%%%%%%%%%%%%%

\documentclass[letterpaper, 10 pt, conference]{ieeeconf}

\IEEEoverridecommandlockouts
\overrideIEEEmargins

\usepackage{amsmath,amssymb}
\usepackage{graphicx}
\usepackage{booktabs}
\usepackage{xcolor}

\newcommand{\todo}[1]{\textcolor{red}{TODO: #1}}
\newcommand{\method}{ToolFrame}

\title{\LARGE \bf
\method{}: Let the Tool Define the Frame for Dexterous Manipulation
}

% RA-L initial submissions follow the double-anonymous review policy.
% Leave the author block without identifying information.
\author{}

\begin{document}

\maketitle
\thispagestyle{empty}
\pagestyle{empty}

\begin{abstract}
The central difficulty in dexterous tool use is that success is determined by
an unobserved functional mechanism rather than by object identity alone: a
brush sweeps through bristle contact, an eraser wipes through a flat patch, a
hammer acts through an impact face, and a screwdriver acts through a tip and
shaft axis. A robot therefore must recover both the tool-side mechanism
(\emph{where} the tool should contact and \emph{which axis} should drive the
effect) and the scene-side mechanism (\emph{where} the effect should start,
move, and terminate). Directly asking a VLM to output robot trajectories hides
this mechanism inside free-form actions and makes evaluation vulnerable to
oracle axes, hidden target references, or task-family shortcuts. This paper
proposes \method{}, a tool-centric image-grounded pipeline that exposes the
mechanism as explicit intermediate references. The VLM is restricted to
selecting sparse public evidence: tool-side functional references
\texttt{tool.contact\_point} and \texttt{tool.motion\_axis}, and scene-side
RGB-D references for start, path, end, and direction. A geometry-constrained
compiler then binds the tool mechanism to the scene mechanism and solves
whole-tool $SE(3)$ subgoals, which are executed by a SimToolReal fixed-goal
dexterous policy. Task contracts evaluate the final physical effect and
attribute failures to functional grounding, scene grounding, subgoal validity,
policy execution, or task effect. \todo{Insert final headline results from
frozen launcher\_summary files and corresponding per-task summary.json files.}
\end{abstract}

\section{INTRODUCTION}

\subsection{Bottom-Level Challenge}
A tool is useful because its geometry induces a task mechanism. The robot does
not merely need to know that an object is a brush, eraser, hammer,
screwdriver, or marker; it must know which part of the object transfers force
or contact to the scene and which motion direction produces the intended
effect. This is especially important for dexterous manipulation, where the hand
may hold the same object in many feasible ways, but only a small subset of
whole-tool poses activates the correct functional part.

\subsection{Mechanistic Gap}
The challenge is to recover this mechanism from public evidence without
collapsing it into an opaque policy or prompt-generated motion. If a VLM is
asked to generate dense waypoints or robot trajectories, the paper can no
longer cleanly answer whether success came from visual grounding, geometric
reasoning, controller robustness, or hidden task priors. The formulation must
therefore expose the causal path from public evidence to tool contact, from
tool contact to scene effect, and from scene effect to task-contract success.

\subsection{Our Solution}
\method{} decomposes the mechanism into two sparse reference sets. The
tool-side set describes how the tool can create an effect: the functional part,
contact point, and motion axis. The scene-side set describes where the effect
should occur: start, path, end, and direction references lifted from RGB-D. A
compiler then solves the missing bridge between them: whole-tool $SE(3)$ poses
that make the tool mechanism coincide with the scene mechanism while respecting
clearance and stage constraints. In short, \emph{the tool defines the frame;
the hand follows}.

\subsection{Contributions}
\begin{enumerate}
  \item A mechanism-explicit, public-evidence tool-use pipeline for dexterous
  manipulation that combines RGB-D scene evidence with public tool-part
  evidence.
  \item A tool-centric reference interface that represents functional tool
  geometry by \texttt{tool.contact\_point} and
  \texttt{tool.motion\_axis}, while restricting the VLM to sparse scene and
  tool references.
  \item A geometry-constrained compiler that mechanistically binds tool refs to
  scene refs and converts them into whole-tool $SE(3)$ subgoals for fixed-goal
  dexterous rollout.
  \item A task-contract benchmark protocol that evaluates the final physical
  effect and attributes failures across grounding, compilation, execution, and
  task-effect stages.
\end{enumerate}

\section{RELATED WORK}

\subsection{Vision-Language Models for Robot Planning}
Discuss VLM/LLM robot planning, visual grounding, language-conditioned
manipulation, and code-as-policies. The underlying challenge is that language
models can encode semantic priors but do not by themselves guarantee metric
contact geometry. The contrast is that \method{} uses VLMs to expose sparse
evidence for the mechanism, not to generate the mechanism's robot trajectory.

\subsection{Tool Affordances and Functional Part Grounding}
Discuss affordance learning, part segmentation, functional region prediction,
and contact-point estimation. The underlying challenge is that a functional
part label is not yet an executable manipulation constraint. The contrast is
that \method{} turns functional parts into explicit contact-point and
motion-axis constraints.

\subsection{Constraint-Based Manipulation and Keyframe Planning}
Discuss task-space constraints, contact-rich manipulation, keyframe/subgoal
planning, and $SE(3)$ optimization. The underlying challenge is to translate a
desired task effect into robot-feasible poses. The contrast is that \method{}
performs this translation at the whole-tool level, binding tool mechanism and
scene mechanism before policy rollout.

\subsection{Tool-Use Benchmarks}
Discuss tool-use and manipulation benchmarks. The contrast is that task
contracts separate the final physical effect from the first failed mechanistic
stage: functional grounding, scene grounding, subgoal validity, policy
execution, or task effect.

\section{PROBLEM FORMULATION AND EVIDENCE BOUNDARY}

\subsection{Mechanistic View of Tool Use}
We formulate a tool-use task as a relation between a tool mechanism and a scene
mechanism. The tool mechanism specifies a contact point and a motion axis in
the tool frame. The scene mechanism specifies the spatial locus and direction
of the desired effect in the world frame. The robot succeeds only if the
executed whole-tool motion makes these two mechanisms coincide under the task
contract.

\subsection{Allowed Public Inputs}
The method may use natural task language, public RGB-D scene evidence, camera
calibration, public tool meshes, segmented tool-part evidence, reviewed or
review-pending public tool references, and robot state needed for execution.

\subsection{Forbidden Method Inputs}
The method must not use hidden evaluator target references, simulator target
centers for snapping, oracle target axes, oracle trajectories, tool-family
fallback, task-family fallback, primitive fallback, repair-loop fallback, or
seed-specific shortcuts.

\subsection{Reference Interface}
Tool references are
\begin{equation}
  \mathcal{R}_{tool} =
  \{p^{contact}_{tool}, a^{motion}_{tool}\},
\end{equation}
where $p^{contact}_{tool}$ is the functional contact point and
$a^{motion}_{tool}$ is the tool-frame motion axis. Scene references are
\begin{equation}
  \mathcal{R}_{scene} =
  \{p^{start}_{scene}, p^{path}_{scene}, p^{end}_{scene},
  a^{direction}_{scene}\}.
\end{equation}
The compiler solves whole-tool poses $T_{world,tool}\in SE(3)$ that align the
tool contact and motion references with the grounded scene references.

\section{METHOD}

\subsection{Overview}
\begin{figure*}[t]
\centering
\fbox{\parbox{0.92\textwidth}{\centering
Placeholder for Fig. 1: (A) Tool-Centric Prior Extraction;
(B) Scene Key-Points Grounding; (C) Geometry-Constrained Optimizer;
(D) Policy Rollout and Evaluation.}}
\caption{\method{} overview. Public tool-part evidence exposes the tool
mechanism; public RGB-D evidence exposes the scene mechanism; a tool-centric
$SE(3)$ compiler binds both into whole-tool subgoals executed by a SimToolReal
fixed-goal dexterous policy. The VLM selects evidence only and does not output
robot trajectories.}
\label{fig:overview}
\end{figure*}

\subsection{A. Tool-Centric Prior Extraction}
The challenge on the tool side is that visual appearance and object category do
not identify the operational mechanism. A hammer handle and hammer face are
both visible parts, but only one creates the strike effect; a marker body and
tip are both graspable, but only the tip creates the trace. This module
therefore extracts a tool-centric prior from public tool evidence: the
functional part, the contact point \texttt{tool.contact\_point}, and the
motion axis \texttt{tool.motion\_axis}. The output is not a robot action; it
is the tool mechanism that later constrains whole-tool poses.

\subsection{B. Scene Key-Points Grounding}
The challenge on the scene side is that the task effect is spatially sparse.
For wiping, sweeping, inserting, or tracing, the important evidence is not a
dense robot path but a small set of scene references that define where the
effect should begin, how it should progress, and where it should terminate.
\method{} uses the VLM only to select these sparse public RGB-D references:
start point, optional path points, end point, and scene direction. The selected
2D references are lifted to 3D using depth and camera calibration.

\subsection{C. Geometry-Constrained Optimizer}
The central mechanistic gap is that tool refs and scene refs live in different
frames. A contact point in the tool frame cannot directly wipe a world-frame
surface, and a scene direction cannot directly command a dexterous hand. The
compiler closes this gap by solving for whole-tool $SE(3)$ poses that satisfy
contact-point matching, motion-axis alignment, clearance, pre-contact approach,
and stage-wise progression through lift, align, approach, and motion phases.
The output is a sequence of fixed whole-tool subgoals, not low-level actions or
VLM-generated trajectories.

\subsection{D. Policy Rollout and Evaluation}
The remaining challenge is execution: even a valid tool-centric subgoal can
fail if the dexterous policy cannot track it, maintain contact, or produce the
desired scene effect. \method{} therefore separates planning from evaluation.
The SimToolReal fixed-goal policy executes the compiled whole-tool subgoals,
while task contracts evaluate the resulting physical effect after rollout and
attribute the first failure to the responsible stage. Contract outcomes are not
used as a repair loop or fallback during the method.

\section{BENCHMARK AND TASKS}

\subsection{Functional Tool Geometry}
\begin{table}[t]
\caption{Functional tool geometries targeted by \method{}.}
\label{tab:tools}
\begin{center}
\begin{tabular}{ll}
\toprule
Tool & Functional mechanism \\
\midrule
Brush & broad lateral contact and sweep axis \\
Eraser & flat bottom contact patch and wipe axis \\
Hammer & impact face and strike axis \\
Screwdriver & tip and shaft/insertion axis \\
Marker & point tip and writing/contact axis \\
\bottomrule
\end{tabular}
\end{center}
\end{table}

\subsection{Scene Coverage}
The benchmark should stress whether the mechanism is genuinely geometric rather
than scene-specific. Tabletop scenes test horizontal manipulation, while
side-board scenes test whether the same tool-centric formulation transfers to
vertical interaction where gravity, camera viewpoint, and contact orientation
change.

\subsection{Task Tiers}
Headline tasks are those with fixed/current contracts and frozen statistics.
Current headline candidates are brush sweep forward and tabletop eraser wipe.
Candidate/provisional tasks include brush left, hammer down, screwdriver wide
socket, and marker trace variants. Diagnostic tasks include harder variants
such as screwdriver deep socket or dense layouts. Draft/smoke tasks are used
for contract development only and must not be reported as final headline
statistics.

\subsection{Metrics}
The metrics diagnose the causal chain from mechanism recovery to task effect:
functional grounding rate, scene grounding rate, subgoal validity rate, policy
execution rate, final task success rate, and failure attribution. Final numbers
must come only from frozen \texttt{launcher\_summary} files and matched
per-task summaries.

\section{EXPERIMENTS}

\subsection{Main Results}
\begin{table*}[t]
\caption{Main benchmark results. All numbers are placeholders until frozen
launcher summaries are available.}
\label{tab:main_results}
\begin{center}
\begin{tabular}{lcccccc}
\toprule
Task & $N$ & Func. & Scene & Subgoal & Policy & Success \\
\midrule
Brush sweep / forward & \todo{N} & \todo{} & \todo{} & \todo{} & \todo{} & \todo{} \\
Eraser wipe / tabletop & \todo{N} & \todo{} & \todo{} & \todo{} & \todo{} & \todo{} \\
Hammer strike / down & \todo{N} & \todo{} & \todo{} & \todo{} & \todo{} & \todo{} \\
Screwdriver insert / wide socket & \todo{N} & \todo{} & \todo{} & \todo{} & \todo{} & \todo{} \\
Marker trace / L or C & \todo{N} & \todo{} & \todo{} & \todo{} & \todo{} & \todo{} \\
\bottomrule
\end{tabular}
\end{center}
\end{table*}

\subsection{Qualitative Rollouts}
Show representative rollouts for each tool family. Each example should expose
the full mechanism chain: functional tool refs, grounded scene refs, compiled
whole-tool subgoals, dexterous rollout, and final task-contract outcome.

\subsection{Failure Attribution}
Use stacked bars or a compact table to show where the mechanism breaks:
functional tool grounding, scene grounding, invalid subgoals, policy execution,
or task effect. This turns failed rollouts into evidence about the pipeline
rather than only negative examples.

\section{FIGURES AND TABLES PLAN}

\begin{itemize}
  \item Fig.~1: challenge-to-solution overview, showing tool mechanism, scene
  mechanism, compiler bridge, and policy rollout.
  \item Fig.~2: five tools and their functional mechanisms: contact points and
  motion axes.
  \item Fig.~3: tabletop and side-board scene mechanisms.
  \item Fig.~4: whole-tool $SE(3)$ subgoals and dexterous rollout.
  \item Fig.~5: failure attribution along the mechanism chain.
  \item Table~I: functional geometry coverage.
  \item Table~II: task, scene, and contract status.
  \item Table~III: main results from frozen summaries.
\end{itemize}

\section{LIMITATIONS}

Mechanism recovery requires visible RGB-D evidence and reasonably accurate
2D-to-3D lifting. Tool-side mechanisms depend on public tool-part evidence and
human review; the current system does not infer arbitrary unseen affordances. A
geometrically valid mechanism alignment can still fail under the fixed-goal
dexterous policy. Side-board and marker-trace tasks may require stronger task
contracts before becoming headline statistics. The system intentionally does
not use online repair, primitive fallback, or seed-specific shortcuts.

\section{CONCLUSIONS}

\method{} treats dexterous tool use as the recovery and execution of a
functional mechanism. The tool supplies the contact and motion frame, the scene
supplies the desired effect frame, and the compiler binds the two into
whole-tool $SE(3)$ subgoals for fixed-goal dexterous execution. By keeping the
VLM restricted to public evidence selection and by evaluating with task
contracts, the paper makes the causal chain from visual grounding to physical
effect explicit. The central message is simple: the tool defines the frame, and
the hand follows.

\bibliographystyle{IEEEtran}
\bibliography{references}

\end{document}
